\documentclass{situationist-preprint}

\title{Public Time and the Datability of Death}
\msid{PHE-07-b}
\msfamily{PHE-07}
\msdiscipline{Phenomenology}
\mstype{Companion manuscript}
\mskeywords{being-towards-death; public time; datability; anticipation; tenseless theory}

\begin{document}
\maketitle

\begin{abstract}
PHE-07-a argues that authentic being-towards-death is possible as an orientation but not as a stable achievement, because anticipation is constantly drawn back into the public time of the \enquote{they}. We examine the mechanism of this drawing back. Public time is datable: every now is a now-when, and in it death appears as an event at some date, happening to someone. Dating one's death converts its indefiniteness into a date and its mineness into an instance of a general fact. We argue that datability is a structure of public time and not a metaphysical present, so that the analysis is compatible with a tenseless theory of time, and that the dating of death invites lapse without excluding retrieval.
\end{abstract}

\section{Datability}

In the final chapter of \emph{Being and Time}, Heidegger describes the time of concern, world-time, by four features: datability, spannedness, significance and publicness \citep[\S\S79--81]{Heidegger1927}. Datability is the structure by which every \enquote{now} is a \enquote{now that\dots}, every \enquote{then} a \enquote{then, when\dots}: times are given as the times of something, and are therefore assignable to dates in a shared reckoning. Public time is the same structure as it is shared, measured by clocks and calendars, and available to anyone.

Two readings of datability must be kept apart. On a tensed reading, datability presupposes a present that is metaphysically privileged, from which other times are dated as past and future. On a structural reading, datability is a feature of how times are given in concern: each time is given as related to events and to a public reckoning, and the \enquote{now} of the dating is the time of the dating act. Blattner has argued that Heidegger's account of world-time is in the first instance an account of how time is understood in everyday activity, not a metaphysics of time \citep{Blattner1999}. We adopt the structural reading. It requires no moving present and is compatible with a tenseless theory on which times are ordered by earlier and later and tensed judgments have tenseless truth conditions \citep{Mellor1998}, in the same way that PHE-05 treats the living present as a structure of experience rather than a metaphysical present.

\section{How dating draws anticipation back}

Anticipation, as Heidegger describes it, relates to death as ownmost, non-relational, certain and indefinite. Public time gives death otherwise.

\begin{proposition}\label{p:date}
Every dating of one's death presents it as an event in public time, and thereby as definite in principle and as an instance of a general fact.
\end{proposition}

\begin{proof}
To date an event is to assign it a place in the public reckoning, at least approximately (\enquote{in some decades}, \enquote{within the year}). An event so placed is definite in principle, since the reckoning has a determinate place for it whether or not anyone knows which, and it is placed in the same order as everyone's deaths, as one instance among them: \enquote{one dies}, as the \enquote{they} says, \enquote{eventually, but not yet}.
\end{proof}

Proposition~\ref{p:date} shows how the falling of Assumption~2.4 of PHE-07-a operates in the case of death. The everyday understanding does not deny death; it dates it, and dating replaces the indefinite certainty that anticipation holds open with a definite but deferred event. Since concern lives in datable time, and every project is dated, the dating of death is not an occasional failure but the default form that the relation to death takes whenever concern resumes.

On the tenseless theory, there is a fact about the date of my death, and Proposition~\ref{p:date} does not deny it. Anticipation is not a refusal of that fact but a way of relating to death that the fact does not exhaust: death as the possibility that bears on every project now, whatever its date. The indefiniteness that anticipation holds open is the indefiniteness of \emph{when} as it bears on one's projects, not a gap in the order of times.

\begin{objection}
Someone told by a physician that she has months to live knows the approximate date of her death. On this account she cannot anticipate it.
\end{objection}

\begin{reply}
The account does not say that dating excludes anticipation; it says that dating draws anticipation back. Knowledge of a near date is the strongest form of that pull, since every arrangement made becomes a dated arrangement. But the relation that anticipation takes up can still be retrieved after each lapse, as Assumption~2.5 of PHE-07-a states: her death bears on what she does today, and not only on a date. This is why the result of PHE-07-a is an orientation and not an achievement, here as elsewhere.
\end{reply}

\section{Conclusion}

Public time draws anticipation back by dating death. Datability is a structure of the time of concern, not a metaphysical present, and the analysis needs only the structure. Dating converts the indefinite, ownmost possibility into an event at some date in a public order, and because concern lives in that order, the conversion recurs whenever concern resumes. Authentic being-towards-death remains possible as an orientation that must be retrieved against this pull, as PHE-07-a concludes.

\begin{thebibliography}{99}\small
\bibitem[Blattner(1999)]{Blattner1999} Blattner, W.~D. (1999). \emph{Heidegger's Temporal Idealism}. Cambridge University Press.
\bibitem[Heidegger(1927)]{Heidegger1927} Heidegger, M. (1927). Sein und Zeit. \emph{Jahrbuch für Philosophie und phänomenologische Forschung} 8.
\bibitem[Mellor(1998)]{Mellor1998} Mellor, D.~H. (1998). \emph{Real Time II}. Routledge.
\end{thebibliography}

\end{document}
